\documentclass[11pt]{article}
\usepackage[margin=1in]{geometry}
\usepackage[T1]{fontenc}
\usepackage{amsmath,amssymb,amsthm,lmodern,microtype,hyperref}
\hypersetup{colorlinks=true,linkcolor=blue,citecolor=blue,urlcolor=blue,
pdftitle={A cone-angle proof of Fuchs's labelled length-ratio conjecture},pdfauthor={Alec Kriebel}}
\title{A cone-angle proof of Fuchs's labelled\\length-ratio conjecture}
\author{Alec Kriebel\\\small Independent researcher\\
\small \href{https://orcid.org/0009-0001-9320-500X}{ORCID: 0009-0001-9320-500X}}
\date{12 October 2026}
\begin{document}
\maketitle
\begin{abstract}
We prove Fuchs's sine formula for the lengths of parallel first-vertex billiard trajectories in every regular \(n\)-gon, \(n\ge3\), with the sine indexed by his endpoint-angle type. A cone-angle index on the double polygon gives the physical type up to a folding-chart sign. Its sine weight is affine invariant; classical Veech cusp reduction then sends an entire direction to polygon chords. The triangle, square and hexagon are treated by marked lattices. Ferudun's 9 October 2026 preprint already proves this same theorem by explicit uniform twists and descent. This note gives a different proof, distinguishes intrinsic and physical labels, and claims neither a new angular invariant nor exclusive priority.
\end{abstract}
\paragraph{Prior work and chronology.}
Veech established the double-polygon surfaces, affine groups and cusp geometry~\cite{Veech,Finster}. Return-angle indices and hyperelliptic transport of marked separatrices have classical precedents~\cite{EMZ,Boissy}. These mechanisms are not new constructions here. Ferudun's Theorem C gives the exact labelled ratio for \(n\ge4\)~\cite[Sections~2--3]{Ferudun}; \(n=3\) is elementary. His reduced angular index agrees with ours. Our proof instead uses all-affine hyperelliptic commutation and the classical cusp theorem.

The repository candidate was posted on 30 September 2026~\cite{Repository}. It already claimed this theorem and used the present cone/hyperelliptic/cusp strategy, but incorrectly asserted chart-independent physical-label invariance. The folding-chart and reversal repair in this note was made on 12 October 2026 UTC, after Ferudun's 9 October release. We do not backdate the repaired proof, certify independent discovery, or claim first or exclusive resolution. The bounded primary-source comparison in the support archive found no explicit all-\(n\) labelled theorem predating the September candidate; this is not an exhaustive absence certificate.

\section{The problem and source conventions}

Let \(n\ge3\), and let a \emph{short trajectory} in a regular Euclidean \(n\)-gon mean a billiard path from a vertex to the first subsequent vertex, reflecting in side interiors. Its Euclidean length is the sum of its segment lengths. Sides are included as limiting boundary trajectories, as in the source.

Fuchs's parallelism identifies unoriented directions under the polygon's dihedral symmetries: the sum or difference of the angles with a side is a multiple of \(2\pi/n\), modulo \(\pi\). The source's type \(A_k\) has \(k\in\mathbb Z/(n-2)\), represented by \(0,\ldots,n-3\). Set
\[
w_n(k)=\sin\frac{(k+1)\pi}{n}.
\]
Here and below this formula uses that canonical representative; it is not a periodic sine function of an arbitrary integer representative.

\textbf{Theorem.} If two parallel short trajectories have types \(A_k,A_\ell\), then
\[
\frac{L_k}{L_\ell}=\frac{w_n(k)}{w_n(\ell)}.
\tag{1}
\]
This is Conjecture 2.7 on printed p.512 of Fuchs's paper~\cite{Fuchs}. There is no small-length limit in the word \emph{short.}

We use Definition 2.1 together with its immediately preceding reflection-parity convention, not the displayed angle equalities in isolation. Label physical polygon vertices clockwise. At a vertex \(P_j\), measure the interior angle counterclockwise from the ray towards \(P_{j-1}\). Write \(\alpha\in[0,(n-2)\pi/n]\) for that angle of the initial path and \(\pi-\beta\) for that angle of the backwards terminal path, as in Fuchs's Figure 8, so \(2\pi/n\le\beta\le\pi\). If the trajectory has \(N\) segments, its type is equivalently
\[
\begin{array}{ll}
N\text{ odd}:&
s=\dfrac{n(\alpha+\beta)}{2\pi}\in\mathbb Z,\qquad k\equiv n-1-s\pmod{n-2},\\[4pt]
N\text{ even}:&
s=\dfrac{n(\beta-\alpha)}{2\pi}\in\mathbb Z,\qquad k\equiv s-1\pmod{n-2}.
\end{array}
\tag{2}
\]
The inequality clauses in Definition 2.1 select the representative when these formulas wrap around. For example, in the sum case the conditions on \(\beta\) give
\[
\max\{0,(n-2k-2)\pi/n\}\le\alpha
\le\min\{(n-2)\pi/n,\,2(n-k-2)\pi/n\},
\]
which are exactly the two sum branches. In the difference case the two possibilities are \(s=k+1\) and \(s=k+3-n\), giving its two branches. At a polygon side the residues \(0\) and \(n-2\) agree. The sign must retain the stated parity when both a sum and a difference happen to be integral multiples; the source explicitly selects that sign before its definition.

We do not assume Conjectures 2.3, 2.4 or 2.6. Proposition 1.2's unlabelled length spectrum alone would not prove (1).

\section{The double polygon and its angle circles}

Let \(P\) be the polygon and \(Q=-P\) its half-turn. Label their vertices \(P_j,Q_j\), with indices modulo \(n\), clockwise and compatibly with this half-turn. Identify every side of \(P\) with its parallel side of \(Q\) by translation, reversing the endpoint order. Denote this translation surface by \(X_n\). For even \(n\) we still use \textbf{two} polygons; we do not replace this surface by the smaller opposite-side quotient.

The corner cycle beginning at \(P_0\) is
\[
P_0,\ Q_1,\ P_2,\ Q_3,\ldots .
\tag{3}
\]
Each corner contributes \(\delta=(n-2)\pi/n\). Thus, for odd \(n\), all vertices give one point with total angle \(2\pi(n-2)\), and the genus is \((n-1)/2\). For even \(n\), they give two points, each of angle \(2\pi r\), where \(r=(n-2)/2\), and the genus is \((n-2)/2\). The two classes are \(P_{\rm even},Q_{\rm odd}\) and \(P_{\rm odd},Q_{\rm even}\).

Reflection in a side line is its midpoint half-turn composed with reflection in the perpendicular bisector of that side; the latter is a symmetry of the regular polygon. Thus ordinary billiard unfolding, as unlabelled polygon sets, produces translates alternately of \(P\) and \(Q\). Passing each copy to \(X_n\) makes a short trajectory a saddle connection, preserving its length and excluding intermediate marked vertices. The first copy is \(P\); the terminal copy is \(P\) for odd \(N\), \(Q\) for even \(N\). Forgetting the individual reflected labels does not change the physical corner angle. This is the usual double-polygon quotient of the unfolding.

The map
\[
\iota:P\longleftrightarrow Q,\qquad z\longmapsto-z
\tag{4}
\]
is an affine involution with derivative \(-I\). It has exactly the \(n\) glued edge midpoints as regular fixed points; it also fixes the unique vertex point when \(n\) is odd, and exchanges the two vertex points when \(n\) is even. For \(n\ge5\) these counts are \(2g+2\). Riemann--Hurwitz shows that the quotient is a sphere, so \(\iota\) is the hyperelliptic involution.

Every orientation-preserving affine automorphism \(f\) of \(X_n\) commutes with \(\iota\). Indeed, \(f\iota f^{-1}\) has derivative \(-I\), hence is holomorphic in the translation complex structure, and is topologically conjugate to \(\iota\). It is therefore a hyperelliptic involution. The conjugate extends across a zero: with flat coordinate \(z=w^r\), derivative \(-I\) lifts to multiplication of \(w\) by an \(r\)th root of \(-1\). Uniqueness of the hyperelliptic involution on a compact Riemann surface of genus at least two gives equality. The same argument works for an orientation-reversing affine map: the conjugate still has complex-linear derivative \(-I\).

At a vertex point of angle \(2\pi r\), the circle of outgoing germs has an angular coordinate modulo \(2\pi r\). An orientation-preserving real-linear derivative \(A\) induces an increasing angular lift \(F\) with
\[
F(\theta+\pi)=F(\theta)+\pi,\qquad
F(\theta+2\pi)=F(\theta)+2\pi.
\tag{5}
\]
One can check (5) directly from \(A(-v)=-A(v)\), continuity and positive orientation. The lift on the cone can have an additional integral sheet offset; that offset cancels in differences at the same point.

\section{The intrinsic index and the physical label}

Let \(\gamma\) be an oriented saddle connection, with outgoing start germ \(u\) and backwards-pointing end germ \(v\).

If its endpoints coincide, their Euclidean directions differ by \(\pi\), so there is a unique residue \(m\pmod r\) such that the positive angular difference from \(u\) to \(v\) is
\[
\pi+2\pi m\pmod{2\pi r}.
\tag{6}
\]
This includes every connection for odd \(n\).

For even \(n\), let \(e=0\) when the endpoints coincide and \(e=1\) when they are different. In the second case bring \(v\) to the start point by \(\iota\). Its direction now agrees with that of \(u\); define \(m\pmod r\) by
\[
\angle(u,\iota v)=2\pi m\pmod{2\pi r}.
\tag{7}
\]
Equations (5) and the commutation with \(\iota\) prove that \(m\), and also \(e\), are preserved by every orientation-preserving affine automorphism.

\textbf{Intrinsic index and the physical folding chart.} Define the intrinsic cone index \(K\) by
\[
\begin{cases}
K\equiv2m+1\pmod{n-2},&n\text{ odd},\\
K\equiv2m+1-e\pmod{n-2},&n\text{ even}.
\end{cases}
\tag{8}
\]

When the initial unfolded corner is normalized to \(P\), with its physical orientation preserved, the source label is \(k=K\). More generally, if the initial folding chart has orientation sign \(\sigma\in\{1,-1\}\), then \(k\equiv\sigma K\pmod{n-2}\). Source labels and intrinsic indices must therefore be distinguished. Here is the corner-by-corner verification for the orientation-preserving initial \(P\) chart, including the endpoint identification in the even case. Normalize the starting corner to \(P_0\), with lower boundary angle zero. For odd \(n\), the index \(t\) of the terminal corner in (3) is even for \(P\), odd for \(Q\), and its lower-boundary cone angle is \(t\delta\). For even \(n\), use \(P_0\) as the start of its own cone cycle and \(Q_0=\iota P_0\) as the start of the other. In the latter case apply \(\iota\) to identify that cone with the starting one; again the lower-boundary cone angle is \(t\delta\). Terminal \(P\) then has \(t=2a+e\), and terminal \(Q\) has \(t=2a+1-e\).

At a terminal \(P\) corner the backward germ has interior offset \(\pi-\beta\). At a terminal \(Q\) corner reflection reverses the offset, giving
\(\delta-(\pi-\beta)=\beta-2\pi/n\).
Consequently, modulo the full cone angle,
\[
(1-e)\pi+2\pi m=
\begin{cases}
t\delta+\pi-(\alpha+\beta),&\text{terminal }P,\\
t\delta+(\beta-\alpha)-2\pi/n,&\text{terminal }Q.
\end{cases}
\tag{9}
\]
For odd \(n\) take \(e=0\) throughout.

For even \(n\), substituting \(t=2a+e\) in the first line gives
\[
mn=a(n-2)+e(n-1)-s.
\]
Since \(n\equiv2\pmod{n-2}\) and \(s\equiv1-k\), this gives \(2m\equiv e-1+k\). Substituting \(t=2a+1-e\) in the second line gives
\[
mn=a(n-2)+s+e-2,
\]
and \(s\equiv k+1\) gives the same conclusion. Changing the full-cone lift in (9) alters the integer equalities by a multiple of \(nr=n(n-2)/2\), which vanishes modulo \(n-2\) because \(n\) is even. For odd \(n\), a lift changes the equalities by a multiple of \(n(n-2)\); the identical calculation with \(t=2a\) or \(2a+1\) and \(e=0\) gives \(2m\equiv k-1\pmod{n-2}\). This proves \(k=K\) for the initial \(P\) chart.

For an initial \(Q\) chart, let \(H\) be a polygon reflection in a center-to-side-midpoint symmetry axis. Its physical fold from \(Q=-P\) to \(P\) is \(H\circ\iota\), whereas \(\iota\) alone is orientation-preserving. The resulting physical path is the polygon reflection by \(H\) of the path obtained from \(\iota\gamma\). Reflection changes the endpoint angles to
\[
\alpha'=\delta-\alpha,\qquad
\beta'=\pi+2\pi/n-\beta.
\]
In (2), this changes \(s\) to \(n-s\) in the sum case and to \(2-s\) in the difference case, so it changes the physical source label to \(-k\) for either segment parity. Since \(\iota\) preserves the intrinsic index, the physical label for a \(Q\)-starting chart is \(-K\). This proves \(k=\sigma K\) for both chart orientations.

Orientation-preserving affine maps preserve \(K\). Under orientation reversal, (6) changes \(m\) to \(-m-1\), while (7) changes \(m\) to \(-m\). Reversing the orientation of a saddle connection on \(X_n\) gives the same changes, using \(\iota^2=1\) in the different-endpoint case. Thus these operations change \(K\) to \(-K\pmod{n-2}\). A source physical label can accordingly be preserved or negated, depending on the initial folding chart.

In particular, reversal of a physical billiard path has
\[
\alpha'=\pi-\beta,\qquad\beta'=\pi-\alpha,\qquad
k'\equiv(-1)^N k\pmod{n-2}.
\]
For odd \(N\), its initial chart remains \(P\), so the label is negated. For even \(N\), the original endpoint lies in \(Q\); the reversed physical initial chart has the opposite orientation, so its label is preserved. This distinction is necessary even for a pentagon: a four-segment path with \(\alpha=\beta=\pi/2\) has label \(2\pmod3\), retained under physical reversal rather than negated.

Crucially,
\[
w_n(-k)=w_n(k).
\tag{10}
\]
For \(k\ne0\) this is the identity
\(\sin((n-1-k)\pi/n)=\sin((k+1)\pi/n)\); for \(k=0\) it is immediate. Consequently the physical weight equals the intrinsic weight, \(w_n(k)=w_n(K)\). It is preserved under affine maps even when the initial folding chart changes orientation, and under all polygon dihedral isometries. No chart-independent invariance of the physical source label itself is asserted.

\section{Reduction of an entire direction to chords}

We invoke the following classical Veech facts for the double regular polygon, for odd \(n\ge5\) and even \(n\ge8\):

\begin{enumerate}
\item It is a lattice translation surface.
\item Every saddle-connection direction is a parabolic direction.
\item For odd \(n\), its Veech group has one cusp, represented by a polygon symmetry direction. For even \(n\), it has two cusps, represented by the two adjacent symmetry directions, separated by \(\pi/n\).

\end{enumerate}
Therefore an orientation-preserving affine automorphism sends any given saddle-connection direction to one of those model directions. This is a direction statement: the \textbf{same} automorphism acts on every connection in that direction.

These are classical results, not conjectures from Fuchs. Veech's Proposition 2.10, Sections 4--7 and Theorem 5.8 give the parabolic direction, double-polygon and cusp facts~\cite{Veech}. Finster's Section 3.2 describes the smaller opposite-side quotient. Her Remark 3.5 transfers its Veech group to the two-polygon cover using Hubert--Schmidt's Lemma J~\cite{HubertSchmidt}; her p.20 identifies the two cusp representatives~\cite{Finster}. The original Veech papers were retrieved and inspected during the present priority audit; the original candidate's access limitation is historical.

In a model direction, every saddle connection is a side or diagonal lying in one polygon. To see this without a length-spectrum assertion, reflect a polygon in its symmetry axis perpendicular to the direction. A vertex not on the axis has a distinct reflected vertex. The line through those two vertices cuts the convex polygon in exactly their chord; an inward ray therefore reaches that reflected vertex before it can reach any side interior. At a vertex on the axis the parallel line is supporting and there is no interior ray. Boundary rays are sides. Applying this to every corner of \(P\) and \(Q\) exhausts the saddle connections.

The direct chord from \(P_0\) to \(P_j\), \(1\le j\le n-1\), has length
\[
2R\sin(j\pi/n),
\tag{11}
\]
where \(R\) is the circumradius. In the source's angles,
\[
\alpha=(n-1-j)\pi/n,\qquad
\beta=(n+1-j)\pi/n.
\]
It has one segment, so (2) gives \(k\equiv j-1\pmod{n-2}\). This also covers \(j=n-1\), the second side of type \(A_0\). Thus every model connection satisfies
\[
L=2R\,w_n(k).
\tag{12}
\]
Rotating its initial vertex to \(P_0\), or applying \(\iota\) to exchange \(P,Q\), preserves the intrinsic index and the length. Thus every model connection, including those beginning in \(Q\), has \(L=2R\,w_n(K)\).

If \(D f=A\) and two connections have a common unoriented direction represented by a unit vector \(v\), both lengths are multiplied by \(\|Av\|\). Their intrinsic weights are preserved. Transporting them to the model direction gives \(L'/w_n(K)=2R\) for each connection. Since each original physical trajectory starts in \(P\), its source label satisfies \(k=K\); the common derivative factor cancels and proves (1) for that direction.

Finally, Fuchs-parallel directions can be aligned by a polygon dihedral isometry, allowing reversal of the connection. For odd \(n\), angles modulo \(\pi\) also permit the half-step \(\pi/n\); the double polygon has the isometry rotating by \(\pi/n\) and exchanging its polygons. For even \(n\), only steps \(2\pi/n\) are needed. All these isometries preserve lengths and weights. This establishes exactly the source's broader parallelism, not only identical unfolded slopes.

\section{The elementary marked lattices}

For \(n=3\), unfold into the triangular lattice. The first vertex on a ray is its primitive lattice point. The direction symmetries preserve this lattice and primitivity, so parallel short trajectories have equal length. There is only type \(A_0\).

For \(n=4\), the same argument uses the square lattice. Types are \(A_1\) when both primitive coordinates are odd and \(A_0\) otherwise, as in Fuchs Section~2.3. Square symmetries preserve this parity condition, so a parallel family has one weight and one length.

For \(n=6\), use the triangular-lattice basis \((1,0),(-1/2,\sqrt3/2)\). The vertices of the developed hexagons are the integer points with \(p+q\not\equiv2\pmod3\). Let \(v=(p,q)\) be primitive in the full triangular lattice and put \(s=p+q\pmod3\). The first hexagon vertex on its positive ray is \(v\) if \(s=0,1\), and \(2v\) if \(s=2\). The source's explicit type rule in Section~2.3 gives:
\[
\begin{array}{c|c|c|c}
s&\text{endpoint}&\text{type}&w_6\\ \hline
0&v&A_1\text{ or }A_3&\sqrt3/2\\
1&v&A_0&1/2\\
2&2v&A_2&1
\end{array}
\]
The triangular-lattice dihedral symmetries preserve the norm and primitivity and send \(s\) to \(s\) or \(-s\); for example their generators \((p,q)\mapsto(q,p)\) and \((p,q)\mapsto(p-q,p)\) do so. Thus a parallel family with \(s=0\) has constant length and weight. In the other family, \(L/w_6=2\|v\|\) in either residue case. This proves (1) for the hexagon and completes the proof.


\section*{Verification and disclosure}
The support archive contains exact finite controls for the vertex gluings, genera, cone indices, parity-sensitive type dictionary, chord labels, sign-invariance of weights and marked lattices, together with an exact pentagon physical-reversal counterexample. They support the geometric proof; they do not replace its universal arguments or establish the cited classical theorems. Floating-point ray tracing is supplementary evidence, not a validated numerical certificate.

AI tools were used extensively in solving, drafting, source comparison, reproduction and independent adversarial checks. This preprint is unrefereed and has not undergone conventional human peer review. The author is responsible for its contents.

\begin{thebibliography}{9}
\bibitem{Fuchs} D. Fuchs, \emph{Billiard trajectories in regular polygons and geodesics on regular polyhedra}, Arnold Math. J. \textbf{7} (2021), 493--517.
\href{https://doi.org/10.1007/s40598-020-00170-8}{doi:10.1007/s40598-020-00170-8}.
\bibitem{Veech} W. A. Veech, \emph{Teichm\"uller curves in moduli space, Eisenstein series and an application to triangular billiards}, Invent. Math. \textbf{97} (1989), 553--583.
\href{https://doi.org/10.1007/BF01388890}{doi:10.1007/BF01388890}.
\bibitem{Finster} M. Finster, \emph{A series of coverings of the regular \(n\)-gon}, arXiv:1005.4588v3 (2011).
\href{https://arxiv.org/abs/1005.4588}{arXiv:1005.4588}.
\bibitem{HubertSchmidt} P. Hubert and T. A. Schmidt, \emph{Invariants of translation surfaces}, Ann. Inst. Fourier \textbf{51} (2001), 461--495.
\href{https://doi.org/10.5802/aif.1829}{doi:10.5802/aif.1829}.
\bibitem{EMZ} A. Eskin, H. Masur and A. Zorich, \emph{Moduli spaces of abelian differentials: the principal boundary, counting problems, and the Siegel--Veech constants}, Publ. Math. IHES \textbf{97} (2003), 61--179.
\href{https://doi.org/10.1007/s10240-003-0015-1}{doi:10.1007/s10240-003-0015-1}.
\bibitem{Boissy} C. Boissy, \emph{Labeled Rauzy classes and framed translation surfaces}, Ann. Inst. Fourier \textbf{65} (2015), 905--932.
\href{https://doi.org/10.5802/aif.2947}{doi:10.5802/aif.2947}.
\bibitem{Ferudun} A. Ferudun, \emph{On seven conjectures of Fuchs on billiard trajectories in regular polygons and geodesics on the regular dodecahedron}, preprint, version 1.0, 9 October 2026.
\href{https://doi.org/10.5281/zenodo.23267641}{doi:10.5281/zenodo.23267641}.
\bibitem{Repository} A. Kriebel, \emph{Prove the labelled sine ratio for short polygon trajectories (5000006)}, repository PR~152, 30 September 2026.
\href{https://github.com/AlecKriebel/Math/pull/152}{github.com/AlecKriebel/Math/pull/152}.
Original submitted head: \texttt{99ef28d28bf380658400886d1f8c2120a5096d97}.
\end{thebibliography}
\end{document}
